% ---
% titre: Problème d'échelle - distance entre deux villes
% theme: FA
% chapitre: Proportionnalité
% sous_chapitre: Résoudre des problèmes
% niveau: [10e]
% type: EVAL
% tags: [proportionnalite, p-question-TS]
% difficulte: 1
% corrige: true
% ---

\question
Sur une carte à l'échelle 1:\,15\,000, la distance mesurée entre deux villes est de 4,2\,cm.

\begin{parts}
\vspace{20pt}\part[2\half] Quelle est la distance réelle entre ces deux villes ? Donne ta réponse en mètres.

\vspace{10pt}{\footnotesize\raggedleft Espace pour ta démarche et tes calculs\par}
\begin{solutionorgrid}[140pt]

\vspace{5pt}
\def\arraystretch{1.5}
\begin{tabularx}{0.6\textwidth}{l|C|C|}
Distance sur la carte [cm]&1&4,2\\
\hline
Distance réelle [cm]&15\,000&\\
\end{tabularx}
\def\arraystretch{1}

\hfill\textbf{(0,5pt)} pour la construction du tableau

\vspace{15pt}$15\,000 \cdot 4,2 : 1 = 63\,000$\,cm \hfill calcul \textbf{(0,5pt)} et réponse \textbf{(0,5pt)}

\vspace{5pt}
Conversion d'unité de longueur: 63\,000\,cm = $63\,000 \cdot 10^{-2} = 630$\,m \hfill\textbf{(0,5pt)}
\end{solutionorgrid}

\begin{tcolorbox}[enhanced jigsaw, interior hidden, colframe=box, parbox=false]%
Réponse: 
\sol{\vspace{20pt}}{La distance réelle entre les deux villes est de 630\,m. \textbf{(0,5pt)}}
\end{tcolorbox}


\vspace{20pt}\part[2\half] Sur cette même carte, quelle distance en centimètres représente 2,85\,km en réalité?

\vspace{10pt}{\footnotesize\raggedleft Espace pour ta démarche et tes calculs\par}
\begin{solutionorgrid}[140pt]

\vspace{5pt}
Conversion d'unité de longueur: 2,85\,km = $2,85 \cdot 10^{5} = 285\,000$\,cm \hfill\textbf{(0,5pt)}

\vspace{5pt}
\def\arraystretch{1.5}
\begin{tabularx}{0.6\textwidth}{l|C|C|}
Distance sur la carte [cm]&1&\\
\hline
Distance réelle [cm]&15\,000&285\,000\\
\end{tabularx}
\def\arraystretch{1}

\hfill\textbf{(0,5pt)} pour la construction du tableau

\vspace{15pt}$285\,000 \cdot 1 : 15\,000 = 19$\,cm \hfill calcul \textbf{(0,5pt)} et réponse \textbf{(0,5pt)}
\end{solutionorgrid}

\begin{tcolorbox}[enhanced jigsaw, interior hidden, colframe=box, parbox=false]%
Réponse: 
\sol{\vspace{20pt}}{2,85\,km représentent 19\,cm sur la carte. \textbf{(0.5pt)}}
\end{tcolorbox}
\end{parts}


